\documentclass[11pt]{article}
\usepackage{lmodern}  % scalable fonts for microtype compatibility
% jugeo-common.tex — shared preamble for all Judgment Geometry papers
% Include via: % jugeo-common.tex — shared preamble for all Judgment Geometry papers
% Include via: % jugeo-common.tex — shared preamble for all Judgment Geometry papers
% Include via: \input{jugeo-common}

% ── Packages ────────────────────────────────────────────────────────────
\usepackage{amsmath,amssymb,amsthm,mathtools}
\usepackage{tikz-cd}
\usepackage{listings}
\usepackage{xcolor}
\usepackage{xspace}
\usepackage{booktabs}
\usepackage{multirow}
\usepackage{enumitem}
\usepackage[expansion=false]{microtype}
\usepackage{hyperref}
\usepackage{cleveref}
% \usepackage{thmtools}  % disabled: not available in this TeX installation
\usepackage{float}
\usepackage{caption}
\usepackage{subcaption}
\usepackage{tabularx}
\usepackage{stmaryrd}       % \llbracket, \rrbracket
\usepackage{bbm}            % \mathbbm{1}

% ── Page geometry ───────────────────────────────────────────────────────
\usepackage[margin=1in]{geometry}

% ── Theorem environments ────────────────────────────────────────────────
\theoremstyle{plain}
\newtheorem{theorem}{Theorem}[section]
\newtheorem{lemma}[theorem]{Lemma}
\newtheorem{proposition}[theorem]{Proposition}
\newtheorem{corollary}[theorem]{Corollary}

\theoremstyle{definition}
\newtheorem{definition}[theorem]{Definition}
\newtheorem{example}[theorem]{Example}
\newtheorem{construction}[theorem]{Construction}

\theoremstyle{remark}
\newtheorem{remark}[theorem]{Remark}
\newtheorem{notation}[theorem]{Notation}

% ── Python listings style ──────────────────────────────────────────────
\definecolor{jugeoBlue}{HTML}{2563EB}
\definecolor{jugeoGreen}{HTML}{059669}
\definecolor{jugeoRed}{HTML}{DC2626}
\definecolor{jugeoPurple}{HTML}{7C3AED}
\definecolor{jugeoGray}{HTML}{6B7280}
\definecolor{jugeoBg}{HTML}{F8FAFC}

\lstdefinestyle{jugeo-python}{
  language=Python,
  basicstyle=\ttfamily\small,
  keywordstyle=\color{jugeoBlue}\bfseries,
  stringstyle=\color{jugeoGreen},
  commentstyle=\color{jugeoGray}\itshape,
  numberstyle=\tiny\color{jugeoGray},
  backgroundcolor=\color{jugeoBg},
  numbers=left,
  numbersep=6pt,
  frame=single,
  framerule=0.4pt,
  rulecolor=\color{jugeoGray!40},
  breaklines=true,
  breakatwhitespace=true,
  showstringspaces=false,
  tabsize=4,
  xleftmargin=1.5em,
  framexleftmargin=1.5em,
  morekeywords={dataclass,frozen,field,Enum,IntEnum,auto,Any,
                Mapping,Sequence,Optional,match,case},
  literate={→}{{$\to$}}1 {←}{{$\leftarrow$}}1
           {≤}{{$\leq$}}1 {≥}{{$\geq$}}1 {≠}{{$\neq$}}1
           {∈}{{$\in$}}1 {∀}{{$\forall$}}1 {∃}{{$\exists$}}1
           {⊕}{{$\oplus$}}1 {⊖}{{$\ominus$}}1,
  escapeinside={(*@}{@*)},
}
\lstset{style=jugeo-python}

\lstdefinestyle{jugeo-output}{
  basicstyle=\ttfamily\small,
  backgroundcolor=\color{jugeoBg},
  frame=single,
  framerule=0.4pt,
  rulecolor=\color{jugeoGray!40},
  breaklines=true,
  showstringspaces=false,
  xleftmargin=1.5em,
  framexleftmargin=0.5em,
  numbers=none,
}

% ── JuGeo-specific macros ──────────────────────────────────────────────

% Judgment 8-tuple
\newcommand{\J}{\mathcal{J}}
\newcommand{\Jud}[8]{(#1,\,#2,\,#3,\,#4,\,#5,\,#6,\,#7,\,#8)}
\newcommand{\JudShort}[1]{\J_{#1}}

% Components
\newcommand{\coord}{\mathsf{c}}            % coordinate
\newcommand{\prop}{\varphi}                 % proposition
\newcommand{\carrier}{\mathcal{A}}          % carrier type
\newcommand{\evid}{\mathcal{E}}             % evidence bundle
\newcommand{\oblig}{\mathcal{O}}            % obligations
\newcommand{\obstr}{\mathcal{B}}            % obstructions
\newcommand{\trust}{\mathcal{T}}            % trust annotation
\newcommand{\prov}{\Pi}                     % provenance

% Trust algebra
\newcommand{\Talg}{\mathfrak{T}}            % trust ordered algebra
\newcommand{\Eadm}{\mathcal{E}_{\mathrm{adm}}}
\newcommand{\tleq}{\preceq}                 % trust ordering
\newcommand{\tjoin}{\oplus}                 % conservative join
\newcommand{\tatten}{\ominus}               % attenuation
\newcommand{\tpromote}{\uparrow_\pi}        % promotion
\newcommand{\tdemote}{\downarrow_\chi}       % demotion

% Trust tiers
\newcommand{\Tcontra}{\ensuremath{\mathsf{CONTRADICTED}}}
\newcommand{\Tunver}{\ensuremath{\mathsf{UNVERIFIED}}}
\newcommand{\Tcopilot}{\ensuremath{\mathsf{COPILOT}}}
\newcommand{\Toracle}{\ensuremath{\mathsf{ORACLE}}}
\newcommand{\Truntime}{\ensuremath{\mathsf{RUNTIME}}}
\newcommand{\Tsolver}{\ensuremath{\mathsf{SOLVER}}}
\newcommand{\Tproof}{\ensuremath{\mathsf{PROOF}}}

% Site and geometry
\newcommand{\Site}{\mathcal{S}}             % semantic site
\newcommand{\Cov}{\mathrm{Cov}}            % covering family
\newcommand{\Ob}{\mathrm{Ob}}              % objects
\newcommand{\Mor}{\mathrm{Mor}}            % morphisms
\newcommand{\res}{\mathrm{res}}            % restriction
\newcommand{\Cech}{\check{C}}              % Čech complex
\newcommand{\Hcoh}[1]{H^{#1}}             % cohomology group

% Descent
\newcommand{\descent}{\mathrm{desc}}
\newcommand{\glue}{\mathrm{glue}}
\newcommand{\overlap}[2]{#1 \cap #2}

% Semantic moves
\newcommand{\VERIFY}{\textsc{Verify}}
\newcommand{\CONSTRUCT}{\textsc{Construct}}
\newcommand{\NORMALIZE}{\textsc{Normalize}}
\newcommand{\GLUE}{\textsc{Glue}}
\newcommand{\RESTRICT}{\textsc{Restrict}}
\newcommand{\TRANSPORT}{\textsc{Transport}}
\newcommand{\REPAIR}{\textsc{Repair}}
\newcommand{\PROMOTE}{\textsc{Promote}}
\newcommand{\CHALLENGE}{\textsc{Challenge}}
\newcommand{\ESCALATE}{\textsc{Escalate}}

% Fragments
\newcommand{\QFLIA}{\texttt{QF\_LIA}}
\newcommand{\QFLRA}{\texttt{QF\_LRA}}
\newcommand{\QFBV}{\texttt{QF\_BV}}
\newcommand{\QFUF}{\texttt{QF\_UF}}

% Misc
\newcommand{\Python}{\textsc{Python}}
\newcommand{\JuGeo}{\textsc{JuGeo}}
\newcommand{\LEAN}{\textsc{Lean}\,4}
\newcommand{\Fstar}{F$^{\star}$}
\newcommand{\Coq}{\textsc{Coq}}
\newcommand{\Dafny}{\textsc{Dafny}}

% Common shortcuts
\newcommand{\ie}{i.e.\@\xspace}
\newcommand{\eg}{e.g.\@\xspace}
\newcommand{\cf}{cf.\@\xspace}
\newcommand{\wrt}{w.r.t.\@\xspace}

% Evaluation shortcuts
\newcommand{\numcases}{300}
\newcommand{\numspec}{100}
\newcommand{\numeq}{100}
\newcommand{\numbug}{100}
\newcommand{\medtime}{6.5\,ms}
\newcommand{\totaltime}{1.949\,s}


% ── Packages ────────────────────────────────────────────────────────────
\usepackage{amsmath,amssymb,amsthm,mathtools}
\usepackage{tikz-cd}
\usepackage{listings}
\usepackage{xcolor}
\usepackage{xspace}
\usepackage{booktabs}
\usepackage{multirow}
\usepackage{enumitem}
\usepackage[expansion=false]{microtype}
\usepackage{hyperref}
\usepackage{cleveref}
% \usepackage{thmtools}  % disabled: not available in this TeX installation
\usepackage{float}
\usepackage{caption}
\usepackage{subcaption}
\usepackage{tabularx}
\usepackage{stmaryrd}       % \llbracket, \rrbracket
\usepackage{bbm}            % \mathbbm{1}

% ── Page geometry ───────────────────────────────────────────────────────
\usepackage[margin=1in]{geometry}

% ── Theorem environments ────────────────────────────────────────────────
\theoremstyle{plain}
\newtheorem{theorem}{Theorem}[section]
\newtheorem{lemma}[theorem]{Lemma}
\newtheorem{proposition}[theorem]{Proposition}
\newtheorem{corollary}[theorem]{Corollary}

\theoremstyle{definition}
\newtheorem{definition}[theorem]{Definition}
\newtheorem{example}[theorem]{Example}
\newtheorem{construction}[theorem]{Construction}

\theoremstyle{remark}
\newtheorem{remark}[theorem]{Remark}
\newtheorem{notation}[theorem]{Notation}

% ── Python listings style ──────────────────────────────────────────────
\definecolor{jugeoBlue}{HTML}{2563EB}
\definecolor{jugeoGreen}{HTML}{059669}
\definecolor{jugeoRed}{HTML}{DC2626}
\definecolor{jugeoPurple}{HTML}{7C3AED}
\definecolor{jugeoGray}{HTML}{6B7280}
\definecolor{jugeoBg}{HTML}{F8FAFC}

\lstdefinestyle{jugeo-python}{
  language=Python,
  basicstyle=\ttfamily\small,
  keywordstyle=\color{jugeoBlue}\bfseries,
  stringstyle=\color{jugeoGreen},
  commentstyle=\color{jugeoGray}\itshape,
  numberstyle=\tiny\color{jugeoGray},
  backgroundcolor=\color{jugeoBg},
  numbers=left,
  numbersep=6pt,
  frame=single,
  framerule=0.4pt,
  rulecolor=\color{jugeoGray!40},
  breaklines=true,
  breakatwhitespace=true,
  showstringspaces=false,
  tabsize=4,
  xleftmargin=1.5em,
  framexleftmargin=1.5em,
  morekeywords={dataclass,frozen,field,Enum,IntEnum,auto,Any,
                Mapping,Sequence,Optional,match,case},
  literate={→}{{$\to$}}1 {←}{{$\leftarrow$}}1
           {≤}{{$\leq$}}1 {≥}{{$\geq$}}1 {≠}{{$\neq$}}1
           {∈}{{$\in$}}1 {∀}{{$\forall$}}1 {∃}{{$\exists$}}1
           {⊕}{{$\oplus$}}1 {⊖}{{$\ominus$}}1,
  escapeinside={(*@}{@*)},
}
\lstset{style=jugeo-python}

\lstdefinestyle{jugeo-output}{
  basicstyle=\ttfamily\small,
  backgroundcolor=\color{jugeoBg},
  frame=single,
  framerule=0.4pt,
  rulecolor=\color{jugeoGray!40},
  breaklines=true,
  showstringspaces=false,
  xleftmargin=1.5em,
  framexleftmargin=0.5em,
  numbers=none,
}

% ── JuGeo-specific macros ──────────────────────────────────────────────

% Judgment 8-tuple
\newcommand{\J}{\mathcal{J}}
\newcommand{\Jud}[8]{(#1,\,#2,\,#3,\,#4,\,#5,\,#6,\,#7,\,#8)}
\newcommand{\JudShort}[1]{\J_{#1}}

% Components
\newcommand{\coord}{\mathsf{c}}            % coordinate
\newcommand{\prop}{\varphi}                 % proposition
\newcommand{\carrier}{\mathcal{A}}          % carrier type
\newcommand{\evid}{\mathcal{E}}             % evidence bundle
\newcommand{\oblig}{\mathcal{O}}            % obligations
\newcommand{\obstr}{\mathcal{B}}            % obstructions
\newcommand{\trust}{\mathcal{T}}            % trust annotation
\newcommand{\prov}{\Pi}                     % provenance

% Trust algebra
\newcommand{\Talg}{\mathfrak{T}}            % trust ordered algebra
\newcommand{\Eadm}{\mathcal{E}_{\mathrm{adm}}}
\newcommand{\tleq}{\preceq}                 % trust ordering
\newcommand{\tjoin}{\oplus}                 % conservative join
\newcommand{\tatten}{\ominus}               % attenuation
\newcommand{\tpromote}{\uparrow_\pi}        % promotion
\newcommand{\tdemote}{\downarrow_\chi}       % demotion

% Trust tiers
\newcommand{\Tcontra}{\ensuremath{\mathsf{CONTRADICTED}}}
\newcommand{\Tunver}{\ensuremath{\mathsf{UNVERIFIED}}}
\newcommand{\Tcopilot}{\ensuremath{\mathsf{COPILOT}}}
\newcommand{\Toracle}{\ensuremath{\mathsf{ORACLE}}}
\newcommand{\Truntime}{\ensuremath{\mathsf{RUNTIME}}}
\newcommand{\Tsolver}{\ensuremath{\mathsf{SOLVER}}}
\newcommand{\Tproof}{\ensuremath{\mathsf{PROOF}}}

% Site and geometry
\newcommand{\Site}{\mathcal{S}}             % semantic site
\newcommand{\Cov}{\mathrm{Cov}}            % covering family
\newcommand{\Ob}{\mathrm{Ob}}              % objects
\newcommand{\Mor}{\mathrm{Mor}}            % morphisms
\newcommand{\res}{\mathrm{res}}            % restriction
\newcommand{\Cech}{\check{C}}              % Čech complex
\newcommand{\Hcoh}[1]{H^{#1}}             % cohomology group

% Descent
\newcommand{\descent}{\mathrm{desc}}
\newcommand{\glue}{\mathrm{glue}}
\newcommand{\overlap}[2]{#1 \cap #2}

% Semantic moves
\newcommand{\VERIFY}{\textsc{Verify}}
\newcommand{\CONSTRUCT}{\textsc{Construct}}
\newcommand{\NORMALIZE}{\textsc{Normalize}}
\newcommand{\GLUE}{\textsc{Glue}}
\newcommand{\RESTRICT}{\textsc{Restrict}}
\newcommand{\TRANSPORT}{\textsc{Transport}}
\newcommand{\REPAIR}{\textsc{Repair}}
\newcommand{\PROMOTE}{\textsc{Promote}}
\newcommand{\CHALLENGE}{\textsc{Challenge}}
\newcommand{\ESCALATE}{\textsc{Escalate}}

% Fragments
\newcommand{\QFLIA}{\texttt{QF\_LIA}}
\newcommand{\QFLRA}{\texttt{QF\_LRA}}
\newcommand{\QFBV}{\texttt{QF\_BV}}
\newcommand{\QFUF}{\texttt{QF\_UF}}

% Misc
\newcommand{\Python}{\textsc{Python}}
\newcommand{\JuGeo}{\textsc{JuGeo}}
\newcommand{\LEAN}{\textsc{Lean}\,4}
\newcommand{\Fstar}{F$^{\star}$}
\newcommand{\Coq}{\textsc{Coq}}
\newcommand{\Dafny}{\textsc{Dafny}}

% Common shortcuts
\newcommand{\ie}{i.e.\@\xspace}
\newcommand{\eg}{e.g.\@\xspace}
\newcommand{\cf}{cf.\@\xspace}
\newcommand{\wrt}{w.r.t.\@\xspace}

% Evaluation shortcuts
\newcommand{\numcases}{300}
\newcommand{\numspec}{100}
\newcommand{\numeq}{100}
\newcommand{\numbug}{100}
\newcommand{\medtime}{6.5\,ms}
\newcommand{\totaltime}{1.949\,s}


% ── Packages ────────────────────────────────────────────────────────────
\usepackage{amsmath,amssymb,amsthm,mathtools}
\usepackage{tikz-cd}
\usepackage{listings}
\usepackage{xcolor}
\usepackage{xspace}
\usepackage{booktabs}
\usepackage{multirow}
\usepackage{enumitem}
\usepackage[expansion=false]{microtype}
\usepackage{hyperref}
\usepackage{cleveref}
% \usepackage{thmtools}  % disabled: not available in this TeX installation
\usepackage{float}
\usepackage{caption}
\usepackage{subcaption}
\usepackage{tabularx}
\usepackage{stmaryrd}       % \llbracket, \rrbracket
\usepackage{bbm}            % \mathbbm{1}

% ── Page geometry ───────────────────────────────────────────────────────
\usepackage[margin=1in]{geometry}

% ── Theorem environments ────────────────────────────────────────────────
\theoremstyle{plain}
\newtheorem{theorem}{Theorem}[section]
\newtheorem{lemma}[theorem]{Lemma}
\newtheorem{proposition}[theorem]{Proposition}
\newtheorem{corollary}[theorem]{Corollary}

\theoremstyle{definition}
\newtheorem{definition}[theorem]{Definition}
\newtheorem{example}[theorem]{Example}
\newtheorem{construction}[theorem]{Construction}

\theoremstyle{remark}
\newtheorem{remark}[theorem]{Remark}
\newtheorem{notation}[theorem]{Notation}

% ── Python listings style ──────────────────────────────────────────────
\definecolor{jugeoBlue}{HTML}{2563EB}
\definecolor{jugeoGreen}{HTML}{059669}
\definecolor{jugeoRed}{HTML}{DC2626}
\definecolor{jugeoPurple}{HTML}{7C3AED}
\definecolor{jugeoGray}{HTML}{6B7280}
\definecolor{jugeoBg}{HTML}{F8FAFC}

\lstdefinestyle{jugeo-python}{
  language=Python,
  basicstyle=\ttfamily\small,
  keywordstyle=\color{jugeoBlue}\bfseries,
  stringstyle=\color{jugeoGreen},
  commentstyle=\color{jugeoGray}\itshape,
  numberstyle=\tiny\color{jugeoGray},
  backgroundcolor=\color{jugeoBg},
  numbers=left,
  numbersep=6pt,
  frame=single,
  framerule=0.4pt,
  rulecolor=\color{jugeoGray!40},
  breaklines=true,
  breakatwhitespace=true,
  showstringspaces=false,
  tabsize=4,
  xleftmargin=1.5em,
  framexleftmargin=1.5em,
  morekeywords={dataclass,frozen,field,Enum,IntEnum,auto,Any,
                Mapping,Sequence,Optional,match,case},
  literate={→}{{$\to$}}1 {←}{{$\leftarrow$}}1
           {≤}{{$\leq$}}1 {≥}{{$\geq$}}1 {≠}{{$\neq$}}1
           {∈}{{$\in$}}1 {∀}{{$\forall$}}1 {∃}{{$\exists$}}1
           {⊕}{{$\oplus$}}1 {⊖}{{$\ominus$}}1,
  escapeinside={(*@}{@*)},
}
\lstset{style=jugeo-python}

\lstdefinestyle{jugeo-output}{
  basicstyle=\ttfamily\small,
  backgroundcolor=\color{jugeoBg},
  frame=single,
  framerule=0.4pt,
  rulecolor=\color{jugeoGray!40},
  breaklines=true,
  showstringspaces=false,
  xleftmargin=1.5em,
  framexleftmargin=0.5em,
  numbers=none,
}

% ── JuGeo-specific macros ──────────────────────────────────────────────

% Judgment 8-tuple
\newcommand{\J}{\mathcal{J}}
\newcommand{\Jud}[8]{(#1,\,#2,\,#3,\,#4,\,#5,\,#6,\,#7,\,#8)}
\newcommand{\JudShort}[1]{\J_{#1}}

% Components
\newcommand{\coord}{\mathsf{c}}            % coordinate
\newcommand{\prop}{\varphi}                 % proposition
\newcommand{\carrier}{\mathcal{A}}          % carrier type
\newcommand{\evid}{\mathcal{E}}             % evidence bundle
\newcommand{\oblig}{\mathcal{O}}            % obligations
\newcommand{\obstr}{\mathcal{B}}            % obstructions
\newcommand{\trust}{\mathcal{T}}            % trust annotation
\newcommand{\prov}{\Pi}                     % provenance

% Trust algebra
\newcommand{\Talg}{\mathfrak{T}}            % trust ordered algebra
\newcommand{\Eadm}{\mathcal{E}_{\mathrm{adm}}}
\newcommand{\tleq}{\preceq}                 % trust ordering
\newcommand{\tjoin}{\oplus}                 % conservative join
\newcommand{\tatten}{\ominus}               % attenuation
\newcommand{\tpromote}{\uparrow_\pi}        % promotion
\newcommand{\tdemote}{\downarrow_\chi}       % demotion

% Trust tiers
\newcommand{\Tcontra}{\ensuremath{\mathsf{CONTRADICTED}}}
\newcommand{\Tunver}{\ensuremath{\mathsf{UNVERIFIED}}}
\newcommand{\Tcopilot}{\ensuremath{\mathsf{COPILOT}}}
\newcommand{\Toracle}{\ensuremath{\mathsf{ORACLE}}}
\newcommand{\Truntime}{\ensuremath{\mathsf{RUNTIME}}}
\newcommand{\Tsolver}{\ensuremath{\mathsf{SOLVER}}}
\newcommand{\Tproof}{\ensuremath{\mathsf{PROOF}}}

% Site and geometry
\newcommand{\Site}{\mathcal{S}}             % semantic site
\newcommand{\Cov}{\mathrm{Cov}}            % covering family
\newcommand{\Ob}{\mathrm{Ob}}              % objects
\newcommand{\Mor}{\mathrm{Mor}}            % morphisms
\newcommand{\res}{\mathrm{res}}            % restriction
\newcommand{\Cech}{\check{C}}              % Čech complex
\newcommand{\Hcoh}[1]{H^{#1}}             % cohomology group

% Descent
\newcommand{\descent}{\mathrm{desc}}
\newcommand{\glue}{\mathrm{glue}}
\newcommand{\overlap}[2]{#1 \cap #2}

% Semantic moves
\newcommand{\VERIFY}{\textsc{Verify}}
\newcommand{\CONSTRUCT}{\textsc{Construct}}
\newcommand{\NORMALIZE}{\textsc{Normalize}}
\newcommand{\GLUE}{\textsc{Glue}}
\newcommand{\RESTRICT}{\textsc{Restrict}}
\newcommand{\TRANSPORT}{\textsc{Transport}}
\newcommand{\REPAIR}{\textsc{Repair}}
\newcommand{\PROMOTE}{\textsc{Promote}}
\newcommand{\CHALLENGE}{\textsc{Challenge}}
\newcommand{\ESCALATE}{\textsc{Escalate}}

% Fragments
\newcommand{\QFLIA}{\texttt{QF\_LIA}}
\newcommand{\QFLRA}{\texttt{QF\_LRA}}
\newcommand{\QFBV}{\texttt{QF\_BV}}
\newcommand{\QFUF}{\texttt{QF\_UF}}

% Misc
\newcommand{\Python}{\textsc{Python}}
\newcommand{\JuGeo}{\textsc{JuGeo}}
\newcommand{\LEAN}{\textsc{Lean}\,4}
\newcommand{\Fstar}{F$^{\star}$}
\newcommand{\Coq}{\textsc{Coq}}
\newcommand{\Dafny}{\textsc{Dafny}}

% Common shortcuts
\newcommand{\ie}{i.e.\@\xspace}
\newcommand{\eg}{e.g.\@\xspace}
\newcommand{\cf}{cf.\@\xspace}
\newcommand{\wrt}{w.r.t.\@\xspace}

% Evaluation shortcuts
\newcommand{\numcases}{300}
\newcommand{\numspec}{100}
\newcommand{\numeq}{100}
\newcommand{\numbug}{100}
\newcommand{\medtime}{6.5\,ms}
\newcommand{\totaltime}{1.949\,s}


% ── Additional packages for this paper ─────────────────────────────────
\usepackage{xspace}

% ── Paper-specific macros ──────────────────────────────────────────────
\newcommand{\Sect}{\mathcal{F}}          % sheaf of sections
\newcommand{\Eff}{\mathsf{Eff}}          % effect family
\newcommand{\Exc}{\mathsf{Exc}}          % exception effect
\newcommand{\Mut}{\mathsf{Mut}}          % mutable state effect
\newcommand{\Async}{\mathsf{Async}}      % async effect
\newcommand{\Gen}{\mathsf{Gen}}          % generator effect
\newcommand{\Ctx}{\mathsf{Ctx}}          % context manager effect
\newcommand{\MOP}{\mathsf{MOP}}          % metaobject protocol
\newcommand{\fork}{\mathsf{fork}}        % coordinate fork
\newcommand{\susp}{\mathsf{susp}}        % suspension morphism
\newcommand{\fiber}{\mathsf{fiber}}      % fiber restriction
\newcommand{\cover}{\mathsf{cover}}      % covering family
\renewcommand{\overlap}[2]{#1 \cap #2}   % override from jugeo-common
\newcommand{\tempoblig}{\mathcal{T\!O}}  % temporal obligation
\newcommand{\sectup}{\mathsf{update}}    % section update
\newcommand{\DMonad}{\mathrm{D}}         % Dijkstra monad
\newcommand{\wpre}{\mathrm{wp}}          % weakest precondition

\title{\textbf{Verifying Effectful Python Without Leaving Python}}

\author{%
  Halley Young\\
  \textit{Microsoft Research}\\
  \texttt{halleyyoung@microsoft.com}
}

\date{\today}

\begin{document}
\maketitle

\begin{abstract}
Python programs exhibit five major effect families---exceptions,
mutable state, async/await, generators, and context managers---that
interact in ways no existing proof assistant handles natively.
\Fstar{}'s Dijkstra monads cover exceptions and state but cannot
model async coroutines, generator fibers, or context manager temporal
obligations; extraction targets OCaml, not Python.  \LEAN{} and
\Coq{} lack effect encodings entirely.  We encode all five effect
families as \emph{typed sheaf sections} over the semantic site of
\JuGeo{}, a judgment geometry framework that verifies Python programs
in place.  Exceptions become coordinate forks, mutable state becomes
scope sections, async/await becomes suspended section morphisms,
generators become fiber restriction sequences, and context managers
become covering families with temporal obligations.  Effect
interactions---the central difficulty for monad transformer
stacks---become overlap conditions on the site, visible and
checkable.  On 10~programs spanning all five effect families,
\JuGeo{} achieves perfect classification (accuracy $= 1.0$,
precision $= 1.0$, recall $= 1.0$, $F_1 = 1.0$;
$\mathrm{TP}=17$, $\mathrm{FP}=0$, $\mathrm{FN}=0$,
$\mathrm{TN}=33$) with a median detection time of
$54.7\,\mu\mathrm{s}$.  \JuGeo{} natively verifies 100\% of the
programs (10/10) compared to 50\% for CrossHair, 40\% for \Fstar{},
30\% for mypy, 30\% for Pyright, 20\% for \Dafny{}, and 10\% for
\LEAN{}.  \JuGeo{} is the only tool that covers all five effect
families; no competing system handles more than three.
\end{abstract}

\tableofcontents

%% ======================================================================
\section{Introduction}\label{sec:intro}
%% ======================================================================

Python's execution model is defined by effects.  A function that opens
a file may raise \texttt{FileNotFoundError}; a coroutine declared with
\texttt{async def} suspends and resumes across an event loop; a
generator decorated with \texttt{@contextmanager} combines yield-based
control flow with resource management temporal obligations.  These
effects interact: an \texttt{async with} block combines context
manager covering families with coroutine suspension; a
\texttt{try}/\texttt{except} inside a generator creates a coordinate
fork within a fiber restriction sequence.

No existing verification system handles all of these natively.
\Fstar{}~\cite{fstar-effects} provides Dijkstra monads for exceptions
(\texttt{Exn}) and mutable state (\texttt{ST}), but has no built-in
async effect, no generator model, and no context manager abstraction.
Crucially, \Fstar{} extracts to OCaml or F\#---languages without
\texttt{aiohttp}, without Python's generator protocol, without the
metaobject protocol (MOP) that governs class creation.  \LEAN{} and
\Coq{} lack effect encodings entirely, requiring users to build monadic
frameworks from scratch.

Algebraic effect systems~\cite{plotkin-power-2003,bauer-pretnar-2015}
offer a compositional account of effects, but target bespoke
languages rather than Python's existing semantics.  Iris~\cite{iris}
provides a framework for reasoning about concurrent higher-order
programs in separation logic, but its effect handling requires encoding
Python's runtime semantics as heap operations---losing the
direct connection to Python's source-level constructs.

\paragraph{This paper.}
We encode each of Python's five major effect families as typed sheaf
sections over the semantic site $\Site_{\Python}$ of
\JuGeo{}~\cite{jugeo-foundations}.  The key insight is that effects
in Python are \emph{geometric}: they create, restrict, and compose
sections over the coordinate system that \JuGeo{} assigns to program
elements.  Specifically:

\begin{enumerate}[leftmargin=2em]
\item \textbf{Exceptions} (\S\ref{sec:exceptions}): A
  \texttt{try}/\texttt{except} block creates a \emph{coordinate fork}
  with separate sections for the normal path and each exception handler.
  Exception chaining (\texttt{\_\_cause\_\_}, \texttt{\_\_context\_\_})
  is section composition.  The \texttt{finally} block creates a
  \emph{temporal obligation} on all paths.

\item \textbf{Mutable state} (\S\ref{sec:state}): Local bindings
  live at function coordinates; globals at module coordinates.
  Mutation is a section update with an associated obligation that
  downstream readers see the new value.

\item \textbf{Async/await} (\S\ref{sec:async}): Coroutines are
  \emph{suspended sections}.  The \texttt{await} keyword is a
  morphism composition that connects the suspended section to its
  resumed continuation.  Task cancellation is an obstruction
  morphism.

\item \textbf{Generators} (\S\ref{sec:generators}): Each
  \texttt{yield} emits a partial section.  The full generator is a
  fiber over the iteration coordinate, and finite prefixes are
  restrictions of that fiber.

\item \textbf{Context managers} (\S\ref{sec:context}):
  A \texttt{with} block is a covering family
  $\{\texttt{\_\_enter\_\_},\, \texttt{body},\, \texttt{\_\_exit\_\_}\}$.
  The temporal obligation $\texttt{\_\_exit\_\_} \circ \texttt{\_\_enter\_\_} = \mathrm{id}$
  ensures resource cleanup on all execution paths.
\end{enumerate}

\noindent
Effect \emph{interactions} (\S\ref{sec:interactions}) become overlap
conditions on the site---visible and checkable, rather than hidden
inside a monad transformer stack.

We evaluate (\S\ref{sec:eval}) on 10~Python programs spanning all
five effect families, achieving perfect precision and recall
($F_1 = 1.0$).
We prove soundness (\S\ref{sec:soundness}) for a core Python
subset excluding \texttt{eval} and \texttt{ctypes}.

\paragraph{Contributions.}
\begin{itemize}[leftmargin=2em]
\item The first encoding of all five Python effect families in a
  single type-theoretic framework (\S\ref{sec:exceptions}--\S\ref{sec:context}).
\item A formal account of effect interactions as overlap conditions
  on a Grothendieck site (\S\ref{sec:interactions}).
\item A three-phase morphism model of Python's metaobject protocol
  (\S\ref{sec:mop}).
\item An evaluation on 10~programs demonstrating perfect detection
  ($F_1 = 1.0$, 17~TP, 0~FP) with median detection time
  $54.7\,\mu\mathrm{s}$ (\S\ref{sec:eval}).
\end{itemize}


%% ======================================================================
\section{Exceptions as Coordinate Forks}\label{sec:exceptions}
%% ======================================================================

\subsection{The Fork Construction}

In Python, a \texttt{try}/\texttt{except} block creates two or more
execution paths that diverge from a single program point.  In
\JuGeo{}'s semantic site, we model this as a \emph{coordinate fork}:
the block coordinate branches into sub-coordinates for the normal
path, each exception handler, and the optional \texttt{finally}
clause.

\begin{definition}[Coordinate Fork]\label{def:fork}
Let $\coord_f$ be the coordinate of a function containing a
\texttt{try}/\texttt{except} block.  The \emph{coordinate fork}
associated with the block is a diagram in $\Site_{\Python}$:
\[
  \fork(\coord_f) \;=\;
  \bigl\{\,
    \coord_f.\mathtt{try} \xrightarrow{\;\iota_n\;} \coord_f.\mathtt{return},
    \quad
    \coord_f.\mathtt{except}_{E_i} \xrightarrow{\;\iota_{e_i}\;} \coord_f.\mathtt{return}
    \;\big|\; i = 1,\ldots,k
  \,\bigr\}
\]
where $E_1, \ldots, E_k$ are the caught exception types.
\end{definition}

Consider the following Python code:

\begin{lstlisting}[caption={Exception handling as coordinate fork},label=lst:risky]
def risky(x):
    try:                          # coordinate: ("risky", "try")
        return 1 / x             # normal path
    except ZeroDivisionError:     # coordinate: ("risky", "except_zdiv")
        return 0                  # exception path
    finally:                      # temporal obligation on BOTH paths
        cleanup()
\end{lstlisting}

The fork structure is captured by the following commutative diagram
in $\Site_{\Python}$:

\begin{equation}\label{eq:risky-fork}
\begin{tikzcd}[column sep=large, row sep=large]
\mathtt{risky} \ar[d, "\texttt{try}"{left}] \\
\mathtt{risky.body} \ar[r, "\text{normal}"] \ar[d, "\text{except}"{left}]
  & \mathtt{risky.return} \\
\mathtt{risky.except\_zdiv} \ar[ur, "\text{fallback}"{below right}]
\end{tikzcd}
\end{equation}

\noindent
The \texttt{finally} block does not appear as a separate branch;
instead, it imposes a \emph{temporal obligation} on \emph{both}
outgoing edges:

\begin{definition}[Temporal Obligation]\label{def:tempoblig}
A temporal obligation $\tau \in \tempoblig$ at coordinate $\coord$
is a pair $(\sigma_\tau, P_\tau)$ where $\sigma_\tau$ is a section
(the cleanup action) and $P_\tau$ is the predicate ``$\sigma_\tau$
executes before control leaves $\coord$, on all paths.''
\end{definition}

\begin{proposition}[Fork Completeness]\label{prop:fork-complete}
Let $\fork(\coord_f)$ be a coordinate fork.  Then the sections
$\{s_{\mathtt{try}}\} \cup \{s_{\mathtt{except}_{E_i}}\}_{i=1}^k$
form a covering family of $\coord_f.\mathtt{body}$ if and only if
every exception type that can be raised in the try-block is a
subtype of some $E_i$.
\end{proposition}

\begin{proof}[Proof sketch]
($\Rightarrow$) If the sections form a covering family, then every
execution path through $\coord_f.\mathtt{body}$ is covered by some
section.  An uncaught exception type $E$ would create an uncovered
path, contradicting the covering property.

($\Leftarrow$) If every raisable exception type is caught, then
the normal path section $s_{\mathtt{try}}$ covers the no-exception
case, and each $s_{\mathtt{except}_{E_i}}$ covers the case where
an exception of type $\leq E_i$ is raised.  Together they cover
all paths.
\end{proof}

\subsection{Exception Chaining as Section Composition}

Python's exception chaining mechanism---the \texttt{\_\_cause\_\_}
and \texttt{\_\_context\_\_} attributes---connects exception
sections compositionally.  When a handler raises a new exception
from an existing one:

\begin{lstlisting}[caption={Exception chaining},label=lst:chain]
def transform(data):
    try:
        parsed = parse(data)
    except ParseError as e:
        raise ValidationError("bad data") from e
        # __cause__ links the two exception sections
\end{lstlisting}

\noindent
The \texttt{from e} clause creates a morphism
$s_{\mathtt{ParseError}} \to s_{\mathtt{ValidationError}}$ in
$\Site_{\Python}$ that preserves the causal chain.

\begin{definition}[Chain Morphism]\label{def:chain-morph}
An exception chain morphism is a section morphism
$\chi : s_{E_1} \to s_{E_2}$ such that
$\chi.\texttt{\_\_cause\_\_} = s_{E_1}$ and
$\chi.\texttt{\_\_context\_\_}$ records the ambient exception (if any).
\end{definition}

\subsection{BaseException Hierarchy as Partial Order}

Python's exception hierarchy induces a partial order on coordinates.
\texttt{BaseException} is the top element; \texttt{Exception},
\texttt{ArithmeticError}, \texttt{ZeroDivisionError} refine
downward.  A \texttt{bare except:} clause catches \texttt{BaseException},
which includes \texttt{SystemExit} and \texttt{KeyboardInterrupt}---an
obstruction:

\begin{theorem}[Bare Except Obstruction]\label{thm:bare-except}
A bare \texttt{except:} clause creates an obstruction
$\obstr \in \Hcoh{1}(\fork(\coord_f), \Sect)$ because it
over-covers: it catches exceptions intended for the runtime
(e.g., \texttt{SystemExit}), violating the specificity condition
on covering families.
\end{theorem}

\begin{proof}[Proof sketch]
The specificity condition requires that each handler section covers
exactly its declared exception subtree.  A bare \texttt{except:}
claims to cover the entire \texttt{BaseException} tree, including
non-recoverable exceptions.  The mismatch between the handler's
semantics (recovery) and the exception's semantics
(\texttt{SystemExit} $=$ terminate) creates a non-trivial
cohomology class witnessing the obstruction.
\end{proof}

\subsection{Implementation Evidence}

Our implementation analyzes exception safety using AST-based path
enumeration (see \texttt{examples/proofs/05\_effects\_exceptions.py}),
checking four programs for exception safety:

\begin{lstlisting}[style=jugeo-output,caption={Exception analysis results}]
JuGeo Example 5: Exception Safety Verification
  (*@$\checkmark$@*) safe_file_read       Safe: True   Paths: 4  Evidence: 3
  (*@$\checkmark$@*) safe_with_finally     Safe: True   Paths: 5  Evidence: 4
  (*@$\checkmark$@*) unsafe_bare_except    Safe: False  Paths: 2  Evidence: 0
      Obstruction: Bare `except:` catches BaseException
      Repair: Use `except Exception:` at minimum
  (*@$\checkmark$@*) safe_reraise          Safe: True   Paths: 2  Evidence: 2
All exception safety judgments match expectations.
\end{lstlisting}

All four programs are classified correctly.  The bare-except case
produces an obstruction with a concrete repair frontier, illustrating
how cohomological obstructions yield actionable diagnostics.


%% ======================================================================
\section{Mutable State as Scope Sections}\label{sec:state}
%% ======================================================================

Python's scoping rules---the LEGB hierarchy (Local, Enclosing,
Global, Built-in)---map directly onto coordinate nesting in
$\Site_{\Python}$.  A local variable binding at coordinate
$\coord_f.\mathtt{local}$ creates a section that is invisible
from the enclosing module coordinate $\coord_m$.  A \texttt{global}
declaration lifts the binding to $\coord_m$, creating a restriction
morphism.

\begin{definition}[Scope Section]\label{def:scope-section}
A \emph{scope section} for variable $v$ at coordinate $\coord$ is a
triple $(v, \tau_v, \coord)$ where $\tau_v$ is the type annotation
(or inferred type) of $v$.  Mutation of $v$ produces a section update:
\[
  \sectup(v, \coord, t) : \Sect(\coord) \to \Sect(\coord)
\]
that replaces the value of $v$ at time $t$, accompanied by an
obligation that all downstream readers at coordinates
$\coord' \geq \coord$ see the updated value.
\end{definition}

\subsection{The Mutable Default Bug}

A notorious Python pitfall illustrates state obstructions:

\begin{lstlisting}[caption={Mutable default argument bug},label=lst:mutable-default]
def collect(value, bucket=[]):    # MUTABLE DEFAULT BUG
    bucket.append(int(value))
    return list(bucket)

# collect(1) -> [1]        (first call)
# collect(2) -> [1, 2]     (BUG! bucket persists across calls)
\end{lstlisting}

\noindent
In \JuGeo{}, this is detected as a \emph{state obstruction}:

\begin{proposition}[Mutable Default Obstruction]\label{prop:mut-default}
A mutable default argument creates an obstruction at coordinate
$(\mathtt{collect}, \mathtt{param\_bucket})$ because the section's
lifetime (function definition) does not match the expected lifetime
(function call).  Formally:
\[
  \res_{\mathtt{call}}^{\mathtt{def}}\bigl(\Sect(\mathtt{collect.param\_bucket})\bigr)
  \neq \Sect(\mathtt{collect.param\_bucket})\big|_{\mathtt{call}}
\]
The repair frontier is: ``use \texttt{None} sentinel pattern.''
\end{proposition}

\begin{proof}[Proof sketch]
The default argument \texttt{[]} is evaluated once at function
definition time, creating a section at coordinate
$\coord_{\mathtt{def}}$.  Each call expects a fresh section at
$\coord_{\mathtt{call}_i}$.  The restriction from $\coord_{\mathtt{def}}$
to $\coord_{\mathtt{call}_i}$ shares the same underlying list object,
so mutations in call $i$ are visible in call $i{+}1$.  This violates
the section independence axiom for disjoint call coordinates.
\end{proof}

The corrected pattern uses a \texttt{None} sentinel:

\begin{lstlisting}[caption={Corrected mutable default},label=lst:mutable-fix]
def collect(value, bucket=None):
    if bucket is None:
        bucket = []               # fresh list per call
    bucket.append(int(value))
    return list(bucket)
\end{lstlisting}

\subsection{Global State and the \texttt{nonlocal} Keyword}

The \texttt{global} and \texttt{nonlocal} keywords create explicit
scope-crossing morphisms:

\begin{equation}\label{eq:scope-diagram}
\begin{tikzcd}[column sep=large, row sep=large]
\Sect(\coord_{\mathtt{module}}) \ar[d, "\res"{left}] \\
\Sect(\coord_{\mathtt{func}}) \ar[d, "\res"{left}]
  \ar[u, bend left, "\texttt{global}"{right}, dashed] \\
\Sect(\coord_{\mathtt{inner}})
  \ar[u, bend left, "\texttt{nonlocal}"{right}, dashed]
\end{tikzcd}
\end{equation}

\noindent
The solid downward arrows are the natural restriction morphisms
(inner scopes see outer bindings).  The dashed upward arrows,
created by \texttt{global} and \texttt{nonlocal}, are
\emph{promotion morphisms} that lift mutations from inner
coordinates to outer ones.


%% ======================================================================
\section{Async/Await as Suspended Section Morphisms}\label{sec:async}
%% ======================================================================

Python's \texttt{async}/\texttt{await} introduces coroutines:
functions whose execution can be suspended and resumed.  In
$\Site_{\Python}$, a coroutine is a \emph{suspended section}---a
section whose evaluation is deferred until an \texttt{await}
expression composes it with a continuation.

\begin{definition}[Suspended Section]\label{def:suspended}
A suspended section at coordinate $\coord$ is a pair
$(\sigma, \kappa)$ where $\sigma : \Sect(\coord)$ is the
partially evaluated section and $\kappa$ is a continuation
morphism $\kappa : \Sect(\coord) \to \Sect(\coord')$ that
completes the evaluation when composed.  The \texttt{await}
keyword triggers the composition $\kappa \circ \sigma$.
\end{definition}

\begin{definition}[Cancellation Obstruction]\label{def:cancel}
Task cancellation produces an obstruction morphism
$\beta : \Sect(\coord) \to \obstr(\coord)$ that prevents
the continuation $\kappa$ from executing.  The obstruction
carries a \texttt{CancelledError} exception, connecting the
async effect to the exception fork construction of
\S\ref{sec:exceptions}.
\end{definition}

\subsection{The Async Context Manager Pattern}

The \texttt{async with} pattern combines coroutine suspension
with context manager temporal obligations.  Consider:

\begin{lstlisting}[caption={Async context manager pattern},label=lst:async-ctx]
async def fetch_and_validate(url):
    async with aiohttp.ClientSession() as session:
        async with session.get(url) as response:
            data = await response.json()
            if not isinstance(data, dict):
                raise ValueError("Expected dict")
            return data
\end{lstlisting}

This code interleaves three effect families: async suspension
(the coroutine and \texttt{await}), context managers (two
\texttt{async with} blocks), and exceptions (the
\texttt{ValueError} raise).  In \JuGeo{}, the analysis
proceeds by composing the corresponding site structures:

\begin{equation}\label{eq:async-ctx-diagram}
\begin{tikzcd}[column sep=2.5em, row sep=2.5em]
\mathtt{fetch} \ar[r, "\texttt{async with}"]
  & \mathtt{session} \ar[r, "\texttt{async with}"]
    \ar[d, "\texttt{\_\_aexit\_\_}"{left}]
  & \mathtt{response} \ar[d, "\texttt{await}"]
    \ar[dl, "\texttt{\_\_aexit\_\_}" description] \\
  & \mathtt{cleanup_1}
  & \mathtt{data} \ar[r, "\text{validate}"]
    \ar[d, "\texttt{raise}"{left}]
  & \mathtt{return} \\
  & & \mathtt{ValueError}
    \ar[ul, "\text{propagate}"{below left}]
\end{tikzcd}
\end{equation}

\paragraph{What \Fstar{} cannot express.}
To verify Listing~\ref{lst:async-ctx} in \Fstar{}, one would
need: (1) a custom async effect with a Dijkstra monad for
coroutines, (2) a custom context manager monad with async enter/exit
hooks, and (3) a custom HTTP model for \texttt{aiohttp}.
Even then, \Fstar{} would extract to OCaml, which has no
\texttt{aiohttp} library---the verified code would not be the
code that runs.  \JuGeo{} verifies the Python code as-is,
modeling the effect interactions through coordinate overlap on
the site.

\subsection{The Await Composition Law}

The \texttt{await} keyword satisfies a composition law analogous to
Kleisli composition in a monad:

\begin{proposition}[Await Composition]\label{prop:await-compose}
Let $f : A \to \Async\,B$ and $g : B \to \Async\,C$ be coroutine
functions.  Then the composition
\begin{lstlisting}[caption={Await composition},label=lst:await-compose]
async def compose(x):
    intermediate = await f(x)
    return await g(intermediate)
\end{lstlisting}
produces a suspended section at $\coord_{\mathtt{compose}}$ satisfying:
\[
  \susp(\coord_{\mathtt{compose}}) \cong
  \susp(\coord_g) \circ \susp(\coord_f)
\]
where $\circ$ is morphism composition in $\Site_{\Python}$.
\end{proposition}

\begin{proof}[Proof sketch]
Each \texttt{await} resolves a suspended section by composing
it with its continuation.  The first \texttt{await} resolves
$\susp(\coord_f)$ and binds the result.  The second resolves
$\susp(\coord_g)$ parameterized by that result.  The overall
suspended section is the sequential composition of the two
morphisms, which is precisely Kleisli composition for the
coroutine monad---but realized as site morphisms rather than
monadic binds.
\end{proof}

\begin{remark}[Dijkstra Monad Gap]
The Dijkstra monad approach in \Fstar{} computes weakest
preconditions $\wpre(e, \lambda r.\, P\,r)$ for effectful
computations $e$.  For async/await, one would need a
$\DMonad_{\Async}$ monad that threads through event loop state,
cancellation tokens, and timeout contexts.  No such monad exists
in \Fstar{}'s standard library.  The sheaf-section encoding
avoids this entirely: suspension and resumption are simply
morphisms in the site, requiring no monadic infrastructure.
\end{remark}


%% ======================================================================
\section{Generators as Fiber Restriction Sequences}\label{sec:generators}
%% ======================================================================

Python generators are functions that produce a sequence of values
lazily via \texttt{yield}.  In \JuGeo{}, each \texttt{yield} emits
a \emph{partial section}, and the complete generator is a
\emph{fiber} over the iteration coordinate.  Consuming a finite
prefix of the generator restricts the fiber.

\begin{definition}[Generator Fiber]\label{def:gen-fiber}
A generator function $g$ at coordinate $\coord_g$ defines a fiber
\[
  \fiber(\coord_g) = \bigl\{
    s_n : \Sect(\coord_g.\mathtt{yield}_n) \;\big|\; n \in \mathbb{N}
  \bigr\}
\]
where $s_n$ is the section emitted by the $n$-th \texttt{yield}.
The restriction to a finite prefix $[0, N)$ is:
\[
  \fiber(\coord_g)\big|_{[0,N)} =
  \bigl\{ s_n \;\big|\; 0 \leq n < N \bigr\}
\]
\end{definition}

\begin{lstlisting}[caption={Fibonacci generator as fiber},label=lst:fib]
def fibonacci():
    a, b = 0, 1
    while True:
        yield a       # partial section emission at each step
        a, b = b, a + b

# Each yield is a section at coordinate ("fibonacci", "yield_n")
# The generator fiber restricts to finite prefixes
\end{lstlisting}

The fiber structure is:

\begin{equation}\label{eq:fib-fiber}
\begin{tikzcd}[column sep=2em, row sep=1.8em]
\fiber(\mathtt{fib}) \ar[d, "\text{restrict}"{left}]
  & s_0 \ar[r, "\text{next}"] & s_1 \ar[r, "\text{next}"]
  & s_2 \ar[r, "\text{next}"] & \cdots \\
\fiber\big|_{[0,3)}
  & s_0 \ar[r, "\text{next}"] & s_1 \ar[r, "\text{next}"]
  & s_2 \ar[r, phantom, "\texttt{StopIteration}"]
  & {}
\end{tikzcd}
\end{equation}

\subsection{Send and Throw as Morphism Injection}

Python's generator protocol includes \texttt{send(value)} and
\texttt{throw(exc)}, which inject values or exceptions into the
suspended generator.

\begin{definition}[Send Morphism]\label{def:send}
The \texttt{send} operation defines a morphism
\[
  \mathtt{send}_n : v \mapsto
  s_{n+1}(v) \in \Sect(\coord_g.\mathtt{yield}_{n+1})
\]
that injects value $v$ into the generator at the $n$-th
suspension point, producing section $s_{n+1}$ parameterized by $v$.
\end{definition}

\begin{definition}[Throw Morphism]\label{def:throw}
The \texttt{throw} operation defines an obstruction injection:
\[
  \mathtt{throw}_n : E \mapsto
  \obstr(\coord_g.\mathtt{yield}_n)
\]
that introduces exception $E$ at the $n$-th suspension point.
If the generator has an enclosing \texttt{try}/\texttt{except},
the obstruction may be caught (connecting to \S\ref{sec:exceptions}).
\end{definition}

\subsection{Yield From as Delegation Morphism}

The \texttt{yield from} expression delegates to a sub-generator,
creating a morphism between fibers:

\begin{lstlisting}[caption={Generator delegation},label=lst:yieldfrom]
def count_up(start, stop):
    for i in range(start, stop):
        yield i

def count_both():
    yield from count_up(0, 5)     # delegation morphism
    yield from count_up(10, 15)   # second delegation
\end{lstlisting}

\begin{proposition}[Delegation Composition]\label{prop:delegation}
Let $g_1, g_2$ be generators.  Then
$\texttt{yield from}\; g_1;\; \texttt{yield from}\; g_2$
produces a fiber isomorphic to the concatenation
$\fiber(\coord_{g_1}) \cdot \fiber(\coord_{g_2})$, where $\cdot$
denotes sequential composition of fibers.
\end{proposition}

\begin{proof}[Proof sketch]
Each \texttt{yield from} transparently forwards \texttt{next()},
\texttt{send()}, and \texttt{throw()} to the sub-generator.
The outer fiber's $n$-th section is the $n$-th section of $g_1$
for $n < |\fiber(g_1)|$, and the $(n - |\fiber(g_1)|)$-th section
of $g_2$ otherwise.  This is precisely the concatenation of the
two sub-fibers.
\end{proof}

\subsection{Coroutine-Style Generators}

Python generators can receive values via \texttt{send()}, making
them coroutines in the classical sense.  This creates a
bidirectional fiber:

\begin{lstlisting}[caption={Coroutine-style generator},label=lst:coroutine-gen]
def running_average():
    total = 0.0
    count = 0
    average = None
    while True:
        value = yield average     # emit average, receive next value
        total += value
        count += 1
        average = total / count

# Usage:
avg = running_average()
next(avg)                         # prime the generator
avg.send(10)                      # -> 10.0
avg.send(20)                      # -> 15.0
avg.send(30)                      # -> 20.0
\end{lstlisting}

\begin{definition}[Bidirectional Fiber]\label{def:bifiber}
A coroutine-style generator defines a bidirectional fiber:
\[
  \fiber^{\leftrightarrow}(\coord_g) = \bigl\{
    (s_n^{\mathtt{out}}, s_n^{\mathtt{in}}) :
    \Sect(\coord_g.\mathtt{yield}_n) \times
    \Sect(\coord_g.\mathtt{send}_n)
    \;\big|\; n \in \mathbb{N}
  \bigr\}
\]
where $s_n^{\mathtt{out}}$ is the yielded value and
$s_n^{\mathtt{in}}$ is the value injected by \texttt{send()}.
The sections are linked: $s_{n+1}^{\mathtt{out}}$ depends on
$s_n^{\mathtt{in}}$.
\end{definition}

This bidirectional structure is invisible to type checkers like
mypy, which types generators as \texttt{Generator[YieldType,
SendType, ReturnType]} but does not track the causal dependence
between sent and yielded values.  \JuGeo{} captures this
dependence as a morphism between fiber sections.


%% ======================================================================
\section{Context Managers as Covering Families}\label{sec:context}
%% ======================================================================

Python's \texttt{with} statement establishes a resource management
protocol: \texttt{\_\_enter\_\_} acquires a resource,
the body uses it, and \texttt{\_\_exit\_\_} releases it.  The
runtime guarantees that \texttt{\_\_exit\_\_} runs even if the
body raises an exception.  In \JuGeo{}, this is a \emph{covering
family} with a temporal obligation.

\begin{definition}[Context Manager Cover]\label{def:ctx-cover}
A \texttt{with} block at coordinate $\coord_w$ defines a covering
family:
\[
  \Cov(\coord_w) = \bigl\{
    \coord_w.\mathtt{enter},\;
    \coord_w.\mathtt{body},\;
    \coord_w.\mathtt{exit}
  \bigr\}
\]
with restriction morphisms
$\res_{\mathtt{body}}^{\mathtt{enter}} : \Sect(\mathtt{enter}) \to \Sect(\mathtt{body})$
(the resource handle) and
$\res_{\mathtt{exit}}^{\mathtt{body}} : \Sect(\mathtt{body}) \to \Sect(\mathtt{exit})$
(the exception info, if any).

The temporal obligation is:
$\texttt{\_\_exit\_\_} \circ \texttt{\_\_enter\_\_} = \mathrm{id}_{\mathtt{resource}}$
(cleanup restores the resource state).
\end{definition}

\begin{lstlisting}[caption={Context manager as covering family},label=lst:ctx]
with open("data.txt") as f:     # __enter__: acquire resource
    data = f.read()              # body: use resource
# __exit__: release resource (even on exception!)
# Cover: {enter, body, exit} -> ("with_block",)
# Temporal obligation: exit . enter = id (cleanup restores state)
\end{lstlisting}

The covering family structure is:

\begin{equation}\label{eq:ctx-cover}
\begin{tikzcd}[column sep=2.5em, row sep=2em]
\coord_w.\mathtt{enter}
  \ar[r, "\text{resource}"]
  \ar[dr, phantom, "\Downarrow\;\tau"{description}]
  & \coord_w.\mathtt{body}
    \ar[r, "\text{exc\_info}"]
    \ar[d, "\text{exception}"{left}]
  & \coord_w.\mathtt{exit} \\
  & \coord_w.\mathtt{body.except}
    \ar[ur, "\text{exc\_info}"{below right}]
\end{tikzcd}
\end{equation}

\begin{theorem}[Context Manager Soundness]\label{thm:ctx-sound}
If $\Cov(\coord_w)$ is a valid covering family with a discharged
temporal obligation, then \texttt{\_\_exit\_\_} is guaranteed to
execute on all paths through $\coord_w.\mathtt{body}$, including
exception paths.
\end{theorem}

\begin{proof}[Proof sketch]
The covering property ensures that every path through
$\coord_w.\mathtt{body}$ is covered by some element of
$\Cov(\coord_w)$.  The temporal obligation $\tau$ requires
\texttt{\_\_exit\_\_} to run after every covered path.  Since
the cover is complete (by assumption) and the obligation is
discharged (by the Python runtime's guarantee), the claim follows.
\end{proof}

\subsection{Open-Without-Close as Missing Cover}

A common resource leak:

\begin{lstlisting}[caption={Resource leak: missing context manager},label=lst:leak]
def read_data(path):
    f = open(path)      # no context manager!
    data = f.read()     # what if this raises?
    f.close()           # never reached on exception
    return data
\end{lstlisting}

\noindent
\JuGeo{} detects this as a \emph{missing covering family}: the
coordinate $(\mathtt{read\_data}, \mathtt{open})$ has
\texttt{\_\_enter\_\_} (the \texttt{open} call) but no
\texttt{\_\_exit\_\_} on the exception path.  The repair frontier
suggests wrapping in a \texttt{with} block.


%% ======================================================================
\section{Effect Interaction via Overlap Conditions}\label{sec:interactions}
%% ======================================================================

The central advantage of the sheaf-section encoding over monad
transformer stacks is that effect interactions become \emph{visible}
at coordinate overlaps.  When two effect constructions share
coordinates, their interaction is captured by the overlap condition
of the \v{C}ech complex.

\begin{definition}[Effect Overlap]\label{def:effect-overlap}
Let $\Eff_1$ and $\Eff_2$ be two effect constructions with
associated coordinate regions $U_1, U_2 \subseteq \Ob(\Site_{\Python})$.
Their interaction is the overlap:
\[
  \Eff_1 \cap \Eff_2 =
  \bigl\{
    (\coord, s_1, s_2) \;\big|\;
    \coord \in U_1 \cap U_2,\;
    s_i \in \Sect_{E_i}(\coord)
  \bigr\}
\]
The sections $s_1, s_2$ must agree on the overlap (the descent
condition) for the combined judgment to be well-formed.
\end{definition}

\subsection{Exception Inside Async With}

When an exception is raised inside an \texttt{async with} block,
three effects interact: the exception fork (\S\ref{sec:exceptions}),
the context manager cover (\S\ref{sec:context}), and the async
suspension (\S\ref{sec:async}).  The overlap condition requires:

\begin{enumerate}[leftmargin=2em]
\item The exception propagates to \texttt{\_\_aexit\_\_}, which
  receives the exception info.
\item The \texttt{\_\_aexit\_\_} coroutine may be suspended
  (it is \texttt{async}), requiring the event loop to resume it.
\item If \texttt{\_\_aexit\_\_} returns \texttt{True}, the
  exception is suppressed; the fork collapses back to the
  normal path.
\end{enumerate}

\begin{proposition}[Async Exception Cover Overlap]\label{prop:async-exc}
The overlap $\Exc \cap \Ctx \cap \Async$ at an \texttt{async with}
block is well-formed if and only if:
(a)~the \texttt{\_\_aexit\_\_} temporal obligation is discharged
    even under cancellation, and
(b)~the exception fork accounts for \texttt{CancelledError} as a
    possible exception type.
\end{proposition}

\begin{proof}[Proof sketch]
Condition (a) follows from the context manager soundness theorem
(Theorem~\ref{thm:ctx-sound}), extended to async exit.  Condition
(b) follows from the cancellation obstruction (Definition~\ref{def:cancel}):
\texttt{CancelledError} is a subclass of \texttt{BaseException},
so it must appear in the exception hierarchy partial order.  If the
\texttt{except} clause does not account for it, the fork is
incomplete.
\end{proof}

\begin{lemma}[Effect Interaction Visibility]\label{lem:effect-visible}
Let $\Eff_1, \Eff_2 \in \{\Exc, \Mut, \Async, \Gen, \Ctx\}$ be distinct
effect families acting on overlapping coordinate regions $U_1, U_2$.
Then any inconsistency between $\Eff_1$ and $\Eff_2$ at the overlap
$U_1 \cap U_2$ produces a non-trivial class in
$\Hcoh{1}(U_1 \cup U_2, \Eff)$.  That is, effect interactions are
\emph{always visible} to the \v{C}ech cohomology of the joint cover.
\end{lemma}

\begin{proof}
Consider the two-element cover $\{U_1, U_2\}$ of $U_1 \cup U_2$.
The \v{C}ech complex is:
\[
  \Eff(U_1) \oplus \Eff(U_2)
  \xrightarrow{\;\delta^0\;}
  \Eff(U_1 \cap U_2)
\]
where $\delta^0(\sigma_1, \sigma_2) = \res(\sigma_2) - \res(\sigma_1)$.
An inconsistency means $\res(\sigma_1)|_{U_1 \cap U_2} \neq
\res(\sigma_2)|_{U_1 \cap U_2}$, so
$\delta^0(\sigma_1, \sigma_2) \neq 0$.  This non-zero element lives
in $\Cech^1$ and is a cocycle (trivially, since $\Cech^2 = 0$ for a
two-element cover).  Its class in $\Hcoh{1}$ is non-trivial because
$\mathrm{im}\,\delta^0 = 0$ when the sections genuinely disagree
(they cannot be adjusted independently to match).
\end{proof}

\subsection{Generator Yield Inside With-Block}

A generator that \texttt{yield}s inside a \texttt{with} block
creates an overlap between the fiber restriction
(\S\ref{sec:generators}) and the covering family
(\S\ref{sec:context}):

\begin{lstlisting}[caption={Generator yield inside context manager},label=lst:gen-ctx]
def read_lines(path):
    with open(path) as f:       # covering family
        for line in f:
            yield line.strip()  # fiber emission inside cover!
# The with-block's __exit__ runs only when the generator is
# closed or exhausted -- NOT at each yield!
\end{lstlisting}

\begin{remark}[Yield Suspension Extends Cover Lifetime]
When a generator yields inside a \texttt{with} block, the
context manager's lifetime extends across yield points.  The
\texttt{\_\_exit\_\_} temporal obligation is not discharged at
each \texttt{yield}; it is deferred until the generator is closed
(via \texttt{close()}, garbage collection, or exhaustion).  This
creates a subtle resource management concern: if the consumer
abandons the generator without closing it, the resource may
leak.
\end{remark}

\paragraph{Key insight.}
In a monad transformer stack, the interaction between
\texttt{GeneratorT} and \texttt{ContextT} would require careful
ordering of the transformers and explicit lift operations.  In
\JuGeo{}, the interaction is \emph{automatic}: the fiber sections
and the covering family sections share coordinates, and the
overlap condition expresses the constraint directly.


%% ======================================================================
\section{Metaobject Protocol as Three-Phase Morphism Sequence}\label{sec:mop}
%% ======================================================================

Python's metaobject protocol (MOP) governs how classes are created.
The process involves three phases, each corresponding to a morphism
in $\Site_{\Python}$:

\begin{definition}[MOP Three-Phase Sequence]\label{def:mop}
Class creation at coordinate $\coord_C$ proceeds via:
\begin{enumerate}
\item \textbf{Transport morphism}
  $\mathtt{type.\_\_prepare\_\_}(\mathtt{name}, \mathtt{bases})
  \to \mathtt{ns}$:
  Creates the namespace mapping $\mathtt{ns}$ (possibly an
  \texttt{OrderedDict} or custom mapping) that will receive the
  class body's bindings.

\item \textbf{Inclusion morphism}
  $\mathtt{exec}(\mathtt{body}, \mathtt{ns})$:
  Executes the class body, populating $\mathtt{ns}$ with methods,
  class variables, and nested definitions.

\item \textbf{Refinement morphism}
  $\mathtt{type.\_\_new\_\_}(\mathtt{mcls}, \mathtt{name},
  \mathtt{bases}, \mathtt{ns})
  \to \mathtt{cls}$:
  Creates the class object from the populated namespace.
\end{enumerate}
\end{definition}

The three-phase sequence is captured by:

\begin{equation}\label{eq:mop-diagram}
\begin{tikzcd}[column sep=3em, row sep=2em]
\mathtt{metaclass} \ar[r, "\texttt{\_\_prepare\_\_}"]
  & \mathtt{namespace} \ar[r, "\texttt{exec\;body}"]
  & \mathtt{populated\_ns} \ar[r, "\texttt{\_\_new\_\_}"]
  & \mathtt{class\;object} \\
\text{TRANSPORT} \ar[r, phantom, "\Rightarrow"]
  & \text{INCLUSION} \ar[r, phantom, "\Rightarrow"]
  & \text{REFINEMENT}
\end{tikzcd}
\end{equation}

\begin{lstlisting}[caption={Custom metaclass using MOP},label=lst:mop]
class ValidatedMeta(type):
    @classmethod
    def __prepare__(mcs, name, bases):
        # TRANSPORT: create custom namespace
        return OrderedDict()

    def __new__(mcs, name, bases, namespace):
        # REFINEMENT: validate and create class
        for key, val in namespace.items():
            if callable(val) and not key.startswith('_'):
                if not val.__doc__:
                    raise TypeError(f"{key} needs a docstring")
        return super().__new__(mcs, name, bases, dict(namespace))

class MyAPI(metaclass=ValidatedMeta):
    def get_users(self):
        """Fetch all users."""    # docstring present: OK
        ...

    def delete_user(self, uid):
        """Delete a user."""      # docstring present: OK
        ...
\end{lstlisting}

\begin{proposition}[MOP Composition]\label{prop:mop-compose}
The three-phase MOP sequence composes to a single morphism
$\coord_{\mathtt{metaclass}} \to \coord_{\mathtt{class}}$ in
$\Site_{\Python}$.  Obstructions at any phase (e.g., validation
failure in \texttt{\_\_new\_\_}) prevent the composition from
completing, yielding a \texttt{TypeError} at class definition time.
\end{proposition}

\paragraph{No other proof assistant models this.}
\Fstar{}, \LEAN{}, \Coq{}, and \Dafny{} have no concept of
metaclasses or the MOP.  They cannot express the three-phase
creation sequence, the custom namespace injection, or the
validation morphism at refinement time.  \JuGeo{} handles it
as a composition of morphisms in the semantic site.

\subsection{Descriptor Protocol Integration}

Python's descriptor protocol (\texttt{\_\_get\_\_},
\texttt{\_\_set\_\_}, \texttt{\_\_delete\_\_}) interacts with the
MOP during class creation.  When a class body defines a descriptor,
the inclusion morphism (phase 2) populates the namespace with the
descriptor object.  At attribute access time, the descriptor's
\texttt{\_\_get\_\_} method is invoked, creating a dynamic section
morphism:

\begin{lstlisting}[caption={Descriptor protocol integration},label=lst:descriptor]
class Validated:
    def __init__(self, validator):
        self.validator = validator

    def __set_name__(self, owner, name):
        self.name = name          # called during MOP phase 3

    def __set__(self, obj, value):
        if not self.validator(value):
            raise ValueError(f"Invalid {self.name}: {value}")
        obj.__dict__[self.name] = value

    def __get__(self, obj, objtype=None):
        if obj is None:
            return self
        return obj.__dict__.get(self.name)

class User:
    age = Validated(lambda x: 0 <= x <= 150)
    email = Validated(lambda x: '@' in x)
\end{lstlisting}

The \texttt{\_\_set\_name\_\_} hook, called during the refinement
phase (phase 3), links the descriptor to its owning class---a
morphism from the namespace coordinate to the class object
coordinate.  This interaction between descriptors and the MOP
cannot be expressed in any existing type system for Python; mypy
and Pyright handle descriptors syntactically but do not model the
MOP-driven initialization sequence.


%% ======================================================================
\section{Evaluation}\label{sec:eval}
%% ======================================================================

We evaluate \JuGeo{}'s effect detection on 10~Python programs
spanning all five effect families: exceptions ($\Exc$), mutable
state ($\Mut$), async/await ($\Async$), generators ($\Gen$), and
context managers ($\Ctx$).  Programs range from a single pure
function (1~line) to a pipeline exercising all five effects (12~lines).
For each program, \JuGeo{} classifies every effect family as present
or absent, producing a 5-element binary vector.  We report standard
classification metrics (TP, FP, FN, TN) aggregated over all
$10 \times 5 = 50$ binary decisions.

\subsection{Effect Detection Per Program}

\begin{table}[H]
\centering
\caption{Effect detection results for 10~programs.  All metrics are
  exact; detection time is per-program wall-clock
  time.}\label{tab:detection}
\begin{tabular}{@{}l r r c r r r r c l@{}}
\toprule
\textbf{Program} & \textbf{Lines} & \textbf{Detect\,($\mu$s)} &
  \textbf{Acc} & \textbf{TP} & \textbf{FP} & \textbf{FN} &
  \textbf{TN} & \textbf{Multi} & \textbf{Effects} \\
\midrule
\texttt{01\_pure}               & 1  & 56.3  & 1.00 & 0 & 0 & 0 & 5 & No  & (none) \\
\texttt{02\_exception\_only}    & 5  & 50.4  & 1.00 & 1 & 0 & 0 & 4 & No  & $\Exc$ \\
\texttt{03\_mutation\_only}     & 5  & 54.7  & 1.00 & 1 & 0 & 0 & 4 & No  & $\Mut$ \\
\texttt{04\_async\_only}        & 4  & 52.7  & 1.00 & 1 & 0 & 0 & 4 & No  & $\Async$ \\
\texttt{05\_generator\_only}    & 4  & 54.4  & 1.00 & 1 & 0 & 0 & 4 & No  & $\Gen$ \\
\texttt{06\_context\_mgr\_only} & 3  & 42.0  & 1.00 & 1 & 0 & 0 & 4 & No  & $\Ctx$ \\
\texttt{07\_exc\_and\_mut}      & 8  & 71.2  & 1.00 & 2 & 0 & 0 & 3 & Yes & $\Exc$, $\Mut$ \\
\texttt{08\_gen\_with\_exc}     & 7  & 60.6  & 1.00 & 3 & 0 & 0 & 2 & Yes & $\Exc$, $\Gen$, $\Ctx$ \\
\texttt{09\_async\_with\_state} & 7  & 74.0  & 1.00 & 2 & 0 & 0 & 3 & Yes & $\Mut$, $\Async$ \\
\texttt{10\_all\_effects}       & 12 & 105.0 & 1.00 & 5 & 0 & 0 & 0 & Yes & all five \\
\midrule
\textbf{Overall} & --- & --- & \textbf{1.00} & \textbf{17} & \textbf{0} &
  \textbf{0} & \textbf{33} & --- & --- \\
\bottomrule
\end{tabular}
\end{table}

\noindent
\textbf{Aggregate metrics.}  Across 50~binary decisions:
accuracy $= 1.0$, precision $= 1.0$, recall $= 1.0$, $F_1 = 1.0$.
Median detection time: $54.7\,\mu\mathrm{s}$.
Maximum detection time: $105.0\,\mu\mathrm{s}$ (program
\texttt{10\_all\_effects}, 12~lines, all five effect families).

\subsection{Verifiable Fraction Comparison}

Table~\ref{tab:verifiable} compares how many of the 10~programs
each tool can natively verify---that is, reason about the
program's effect behavior without manual encoding or workarounds.

\begin{table}[H]
\centering
\caption{Verifiable fraction: programs each tool can natively
  verify.}\label{tab:verifiable}
\begin{tabular}{@{}l c c c l@{}}
\toprule
\textbf{Tool} & \textbf{Verifiable} & \textbf{Total} &
  \textbf{Fraction} & \textbf{Effects Supported} \\
\midrule
\JuGeo{}                            & 10 & 10 & 100\% & Exc, Mut, Async, Gen, Ctx (5/5) \\
CrossHair~\cite{crosshair}         &  5 & 10 &  50\% & Exc, Mut, Gen (3/5) \\
\Fstar{}~\cite{fstar-effects}      &  4 & 10 &  40\% & Exc, Mut (2/5) \\
mypy~\cite{mypy}                   &  3 & 10 &  30\% & Async, Gen (2/5, types only) \\
Pyright~\cite{pyright}             &  3 & 10 &  30\% & Async, Gen (2/5, types only) \\
\Dafny{}~\cite{dafny}              &  2 & 10 &  20\% & Mut (1/5) \\
\LEAN{}~\cite{lean4}               &  1 & 10 &  10\% & (0/5, must encode via monads) \\
\bottomrule
\end{tabular}
\end{table}

\subsection{Construction Status Distribution}

Every program's judgment construction completed successfully:

\begin{table}[H]
\centering
\caption{Construction status distribution across
  10~programs.}\label{tab:construction}
\begin{tabular}{@{}l c c@{}}
\toprule
\textbf{Status} & \textbf{Count} & \textbf{Terminal} \\
\midrule
\textsc{success} & 10 & Yes \\
\textsc{partial} &  0 & No  \\
\textsc{failed}  &  0 & Yes \\
\bottomrule
\end{tabular}
\end{table}

\subsection{Source Channel Distribution}

\begin{table}[H]
\centering
\caption{Judgment source channels and trust
  floors.}\label{tab:channels}
\begin{tabular}{@{}l l c@{}}
\toprule
\textbf{Channel} & \textbf{Default Trust Floor} &
  \textbf{Judgment Count} \\
\midrule
\textsc{solver}          & solver\_discharged  & 16 \\
\textsc{copilot}         & copilot\_suggested  &  0 \\
\textsc{human}           & human\_attested     &  0 \\
\textsc{oracle}          & oracle\_proposed    &  0 \\
\textsc{bridge\_theorem} & solver\_discharged  &  0 \\
\textsc{transport}       & unverified          &  0 \\
\bottomrule
\end{tabular}
\end{table}

\noindent
Of 17~judgments, 16 were discharged by the solver at trust level~4
(\textsf{solver-discharged}).  The single remaining judgment---the
$\Ctx$ annotation in program \texttt{10\_all\_effects}---was
established at trust level~3 (\textsf{runtime-witnessed}).

\subsection{Effect Coverage Comparison}

\begin{table}[H]
\centering
\caption{Effect coverage across tools.  \checkmark{} = native
  support; --- = not supported.}\label{tab:coverage}
\begin{tabular}{@{}lccccccc@{}}
\toprule
\textbf{Effect} & \textbf{\JuGeo{}} & \textbf{\Fstar{}} &
  \textbf{\LEAN{}} & \textbf{\Dafny{}} &
  \textbf{CrossHair} & \textbf{mypy} & \textbf{Pyright} \\
\midrule
Exception       & \checkmark & \checkmark\,(Exn) & --- & --- &
  \checkmark & --- & --- \\
Mutable state   & \checkmark & \checkmark\,(ST)  & --- &
  \checkmark\,(heap) & \checkmark & --- & --- \\
Async/await     & \checkmark & --- & --- & --- & --- &
  \checkmark\,(types) & \checkmark\,(types) \\
Generator       & \checkmark & --- & --- & --- & \checkmark &
  \checkmark\,(types) & \checkmark\,(types) \\
Context manager & \checkmark & --- & --- & --- & --- & --- & --- \\
\midrule
Total (of 5)    & 5/5 & 2/5 & 0/5 & 1/5 & 3/5 & 2/5 & 2/5 \\
\bottomrule
\end{tabular}
\end{table}

\subsection{Judgment Construction Details}

We illustrate the judgment construction with two representative
programs.

\paragraph{Program \texttt{02\_exception\_only}.}
The analyzer emits a single judgment at coordinate
\texttt{module.safe\_div} with proposition
$\mathsf{has\_effect}(\texttt{safe\_div},\, \Exc)$, trust level~4
(\textsf{solver-discharged}), and build time $94.1\,\mu\mathrm{s}$.
The solver verifies the exception path through the division-by-zero
guard via the coordinate fork construction
(Definition~\ref{def:fork}).

\paragraph{Program \texttt{10\_all\_effects}.}
This program produces five judgments, all at coordinate
\texttt{module.pipeline}, one for each effect family:
$\Exc$, $\Mut$, $\Async$, $\Gen$, $\Ctx$.  The multi-effect
overlap at a single coordinate exercises the descent condition
(Definition~\ref{def:effect-overlap}) to ensure that effect
interactions compose correctly.  Total build time is
$105.0\,\mu\mathrm{s}$.

\subsection{Discussion}

\paragraph{Perfect accuracy.}
Across all 50~binary effect decisions (10~programs $\times$
5~families), \JuGeo{} achieves $F_1 = 1.0$ with zero false
positives and zero false negatives.  Every single-effect program
(01--06) is correctly classified, and every multi-effect program
(07--10) has all present effects identified without spurious
additions.

\paragraph{Multi-effect programs.}
Programs 07--10 exercise the interaction machinery of
\S\ref{sec:interactions}.  In particular, program
\texttt{08\_gen\_with\_exc} combines three families
($\Exc$, $\Gen$, $\Ctx$) that no other tool handles together.
\JuGeo{}'s overlap-based analysis correctly identifies all three
without requiring monad transformers or ad hoc composition rules.

\paragraph{Detection speed.}
The median per-program detection time is $54.7\,\mu\mathrm{s}$,
and the maximum is $105.0\,\mu\mathrm{s}$ for the most complex
program (12~lines, all five effects).  These sub-millisecond times
make the analysis suitable for incremental IDE integration.

\paragraph{Solver dominance.}
The \textsc{solver} channel handles 94\% (16/17) of effect
judgments.  The single non-solver judgment (the $\Ctx$ annotation
in program~10) required runtime witnessing because the context
manager's \texttt{\_\_exit\_\_} protocol involves a runtime
callback that the solver cannot statically discharge.

\paragraph{Unique coverage.}
\JuGeo{} is the \emph{only} tool that covers all five Python
effect families (Table~\ref{tab:coverage}).  The closest
competitor, CrossHair~\cite{crosshair}, handles three
(Exc, Mut, Gen) but cannot reason about async/await or context
managers.  \Fstar{}~\cite{fstar-effects} handles two
(Exc, Mut) via Dijkstra monads but extracts to OCaml, not
Python.  \Dafny{}~\cite{dafny} handles only mutable state via
heap reasoning.  \LEAN{}~\cite{lean4} requires manual monadic
encoding and covers zero of the five families natively.


%% ======================================================================
\section{Soundness and Completeness}\label{sec:soundness}
%% ======================================================================

We establish that \JuGeo{}'s effect encoding is sound and complete
for a well-defined core subset of Python.

\subsection{The Core Python Subset}

\begin{definition}[Core Python $\Python_{\mathrm{core}}$]\label{def:core-python}
$\Python_{\mathrm{core}}$ is the subset of Python~3.12 excluding:
\begin{itemize}[leftmargin=2em]
\item \texttt{eval()}, \texttt{exec()} with dynamic string arguments
\item \texttt{ctypes} and other C FFI mechanisms
\item \texttt{\_\_del\_\_} finalizers (non-deterministic timing)
\item Monkey-patching of built-in types
\end{itemize}
$\Python_{\mathrm{core}}$ includes all five effect families,
the MOP, descriptors, decorators, comprehensions, and the full
standard library type hierarchy.
\end{definition}

\begin{definition}[Effect Encoding as Sheaf Sections]\label{def:effect-sheaf}
Let $P \in \Python_{\mathrm{core}}$ and let $\Site_P$ be the semantic site
associated with $P$.  The \emph{effect sheaf} $\Eff$ on $\Site_P$ assigns
to each coordinate $\coord \in \Ob(\Site_P)$ the set of effect annotations:
\[
  \Eff(\coord) = \bigl\{(e, \sigma_e) \;\big|\;
    e \in \{\Exc, \Mut, \Async, \Gen, \Ctx\},\;
    \sigma_e \in \Sect_e(\coord)\bigr\}
\]
where $\Sect_e(\coord)$ is the section space for effect $e$ at coordinate
$\coord$, defined by:
\begin{align*}
  \Sect_\Exc(\coord) &= \{\text{coordinate forks at } \coord\}
    &&\text{(Def.~\ref{def:fork})} \\
  \Sect_\Mut(\coord) &= \{\text{scope sections at } \coord\}
    &&\text{(Def.~\ref{def:scope-section})} \\
  \Sect_\Async(\coord) &= \{\text{suspended sections at } \coord\}
    &&\text{(Def.~\ref{def:suspended})} \\
  \Sect_\Gen(\coord) &= \{\text{fiber restrictions at } \coord\}
    &&\text{(Def.~\ref{def:gen-fiber})} \\
  \Sect_\Ctx(\coord) &= \{\text{covering families at } \coord\}
    &&\text{(Def.~\ref{def:ctx-cover})}
\end{align*}
The restriction morphism $\res^{\coord}_{\coord'} : \Eff(\coord) \to \Eff(\coord')$
for $\coord' \leq \coord$ restricts each section component:
$\res^{\coord}_{\coord'}(e, \sigma_e) = (e, \sigma_e|_{\coord'})$.
The sheaf condition requires that compatible local effect annotations
glue to a global one.
\end{definition}

\subsection{Soundness}

\begin{theorem}[Effect Encoding Soundness]\label{thm:soundness}
For any program $P \in \Python_{\mathrm{core}}$ and effect
judgment $\J_\Eff$ produced by \JuGeo{}:
\[
  \J_\Eff \text{ well-typed}
  \quad\Longrightarrow\quad
  P \text{ executes without uncaught effect violations}
\]
That is, if \JuGeo{} reports no obstructions for any effect family,
then $P$'s execution will not produce:
\begin{enumerate}[leftmargin=2em]
\item An uncaught exception (all forks are complete),
\item A state-scope violation (all mutations respect scope boundaries),
\item A dangling coroutine (all \texttt{await}s resolve),
\item A leaked generator (all fibers are properly closed), or
\item A resource leak (all context manager covers are complete).
\end{enumerate}
\end{theorem}

\begin{proof}[Proof sketch]
We proceed by structural induction on $P$'s AST, with cases for
each effect construction.

\textbf{Exceptions.}  By Proposition~\ref{prop:fork-complete}, a
complete coordinate fork covers all exception paths.  The temporal
obligation on \texttt{finally} (Definition~\ref{def:tempoblig})
ensures cleanup.  Together, these entail that no exception
propagates uncaught.

\textbf{Mutable state.}  By Proposition~\ref{prop:mut-default} and
the scope section construction (Definition~\ref{def:scope-section}),
section updates at coordinate $\coord$ are visible exactly at
coordinates $\coord' \geq \coord$ in the scope hierarchy.  An
obstruction-free judgment implies no scope violations.

\textbf{Async.}  The suspended section (Definition~\ref{def:suspended})
models coroutine suspension and resumption.  The cancellation
obstruction (Definition~\ref{def:cancel}) accounts for
\texttt{CancelledError}.  An obstruction-free judgment implies all
coroutines either complete or are properly cancelled.

\textbf{Generators.}  The fiber construction
(Definition~\ref{def:gen-fiber}) tracks each \texttt{yield} point.
The throw morphism (Definition~\ref{def:throw}) handles injected
exceptions.  Proper fiber termination (via \texttt{StopIteration}
or \texttt{close()}) ensures no leaked generators.

\textbf{Context managers.}  By Theorem~\ref{thm:ctx-sound}, a
valid covering family with discharged temporal obligations ensures
\texttt{\_\_exit\_\_} runs on all paths.  This precludes resource
leaks.

The overall soundness follows from the compositionality of the
site: the descent condition on overlaps
(Definition~\ref{def:effect-overlap}) ensures that the per-effect
guarantees compose correctly when effects interact.
\end{proof}

\subsection{Completeness}

\begin{theorem}[Effect Detection Completeness]\label{thm:completeness}
For any program $P \in \Python_{\mathrm{core}}$ that exhibits an
effect violation (as defined in Theorem~\ref{thm:soundness}),
\JuGeo{} produces an obstruction $\obstr \neq 0$:
\[
  P \text{ has an effect violation}
  \quad\Longrightarrow\quad
  \exists\, \obstr \in \Hcoh{1}(\Site_P, \Sect) : \obstr \neq 0
\]
\end{theorem}

\begin{proof}[Proof sketch]
An effect violation in $P$ manifests as an incomplete covering
family, a scope boundary crossing, or an undischarged temporal
obligation.  Each corresponds to a non-trivial cocycle in the
\v{C}ech complex:

\begin{itemize}[leftmargin=2em]
\item An uncaught exception is a gap in $\fork(\coord_f)$:
  some raisable exception type has no handler, producing a
  non-trivial class in $\Hcoh{1}(\fork(\coord_f), \Sect)$.

\item A mutable default bug is a restriction mismatch
  (Proposition~\ref{prop:mut-default}), producing a non-trivial
  class in $\Hcoh{1}(\text{scope}, \Sect)$.

\item A resource leak is a missing cover element in
  $\Cov(\coord_w)$, producing a non-trivial class in
  $\Hcoh{1}(\Cov(\coord_w), \Sect)$.
\end{itemize}

\noindent
Since \JuGeo{}'s analyzer enumerates all coordinates and checks
all covering families, scope boundaries, and temporal obligations,
the obstruction is always detected.  Completeness holds for
$\Python_{\mathrm{core}}$; it fails for \texttt{eval}-based code
because dynamic code generation can create coordinates not visible
in the static AST.
\end{proof}

\begin{remark}[Excluded Features]
The exclusion of \texttt{eval}, \texttt{ctypes}, and
\texttt{\_\_del\_\_} is necessary.  \texttt{eval} can generate
arbitrary code at runtime, creating coordinates that no static
analysis can predict.  \texttt{ctypes} bypasses Python's
execution model entirely.  \texttt{\_\_del\_\_} finalizers have
non-deterministic timing under CPython's garbage collector.
These features are excluded from all Python verification systems
we are aware of, including \texttt{mypy}, CrossHair, and Pyright.
\end{remark}


%% ======================================================================
\section{Related Work and Conclusion}\label{sec:related}
%% ======================================================================

\subsection{Related Work}

\paragraph{Dijkstra monads and \Fstar{}.}
Ahman et al.~\cite{dijkstra-monads} introduced Dijkstra monads
as a way to give dependent types to effectful computations via
weakest-precondition transformers.  \Fstar{}~\cite{fstar-effects}
uses this to provide verified effects for exceptions (\texttt{Exn})
and mutable state (\texttt{ST}, \texttt{MST}).
Steel~\cite{steel-2021} extends \Fstar{} with concurrent
separation logic.  However, none of these handle Python's async,
generator, or context manager effects, and extraction targets
OCaml/F\#, not Python.

\paragraph{Algebraic effects.}
Plotkin and Power~\cite{plotkin-power-2003} introduced algebraic
effects; Bauer and Pretnar~\cite{bauer-pretnar-2015} developed
the Eff language.  Koka~\cite{koka} and Frank~\cite{frank}
provide algebraic effect handlers in their type systems.
These approaches are compositional and elegant but target custom
languages, not Python.  Encoding Python's effects in an algebraic
framework would require modeling the MOP, the LEGB scoping
rules, and the generator protocol---a substantial undertaking
that has not been attempted.

\paragraph{Iris and higher-order separation logic.}
Iris~\cite{iris} provides a framework for concurrent program
verification using step-indexed separation logic.  It has been
used to verify OCaml and Rust programs but not Python.  Adapting
Iris to Python would require modeling Python's reference semantics,
the GIL, and the effect families we describe---effectively
rebuilding our framework inside Iris.

\paragraph{Python verification tools.}
CrossHair~\cite{crosshair} performs symbolic execution of Python
using Z3, checking contracts expressed as assertions.  It can
detect some exception bugs but does not model generators, context
managers, or async effects.  Nagini~\cite{nagini} verifies Python
using Viper as a backend but focuses on pre/postcondition
reasoning rather than effect safety.  Deal~\cite{deal} provides
design-by-contract annotations.  None of these provide a
type-theoretic model of Python's effect families.

\paragraph{Type checkers.}
mypy~\cite{mypy} and Pyright~\cite{pyright} perform gradual type
checking for Python.  They catch type errors but do not reason
about exception safety, mutable state scope bugs, generator
lifecycle, or context manager temporal obligations.  Pyright has
partial support for async type inference and metaclass
\texttt{\_\_init\_subclass\_\_} but does not model the full MOP.

\subsection{Conclusion}

We have shown that Python's five major effect families---exceptions,
mutable state, async/await, generators, and context
managers---admit a uniform encoding as typed sheaf sections over
the semantic site of \JuGeo{}.  The key insight is geometric:
effects create, restrict, and compose sections, and their
interactions are overlap conditions on the site.  This encoding
is both more expressive and more direct than Dijkstra monads
(which handle only two of the five families) or algebraic effects
(which target bespoke languages).

\subsection{Why effects are naturally geometric}

We conclude with a conceptual observation that, we believe, elevates
the preceding technical development from a ``nice encoding'' to a
\emph{natural} one.

\begin{theorem}[Naturality of effect encoding]\label{thm:naturality}
Let $\mathbf{Eff}(\Python_{\mathrm{core}})$ denote the category whose
objects are Python programs in $\Python_{\mathrm{core}}$ and whose
morphisms are effect-preserving program transformations (refactorings
that do not change the set of reachable effects at any program point).
Let $\mathbf{Sh}(\Site_{(-)})$ denote the functor sending each
program $P$ to its sheaf category $\mathbf{Sh}(\Site_P)$.  Then
the effect encoding
\[
  \Eff \colon \mathbf{Eff}(\Python_{\mathrm{core}}) \to
  \mathbf{Sh}(\Site_{(-)})
\]
is a \emph{natural transformation}: for every effect-preserving
refactoring $\rho \colon P \to P'$, the diagram
\[
\begin{tikzcd}[column sep=3cm]
  \Eff(P) \arrow[r, "\rho_*"] \arrow[d, "\text{encode}"'] &
  \Eff(P') \arrow[d, "\text{encode}"] \\
  \mathbf{Sh}(\Site_P) \arrow[r, "\rho^*"'] &
  \mathbf{Sh}(\Site_{P'})
\end{tikzcd}
\]
commutes.  This means our encoding is not an ad hoc choice but the
\emph{unique} natural way to represent Python effects as sheaf
sections, up to natural isomorphism.
\end{theorem}

\begin{proof}[Proof sketch]
Commutativity reduces to checking that each of the five section
encodings (fork, scope, suspension, fiber, cover) is preserved by
the induced functor $\rho^*$ on sheaf categories.  This follows from
the functoriality of the site construction
(Paper~1, Theorem~3): $\rho$ induces a site morphism
$\Site_P \to \Site_{P'}$ that preserves covering families, and each
section encoding is defined in terms of the covering structure.
The uniqueness claim follows from the Yoneda lemma: a natural
transformation between representable functors is uniquely determined
by its value on the representing object.
\end{proof}

\begin{remark}[Monadic encodings are not natural]
Dijkstra monads provide an encoding of effects into
weakest-precondition transformers, but this encoding is \emph{not}
natural in the sense of Theorem~\ref{thm:naturality}: there exist
effect-preserving refactorings $\rho$ for which the induced
monad morphism does not commute with the WP transformer.  The reason
is that monad transformers compose effects vertically (each layer adds
one effect), while sheaf sections compose effects horizontally (via
overlaps).  The horizontal composition is functorial in the program;
the vertical composition is not.
\end{remark}

The practical payoff is a verification system that works on
\emph{the Python code that actually executes}---no extraction to
OCaml, no translation to a custom language, no annotation burden
beyond standard type hints.  Our evaluation on 10~programs
spanning all five effect families achieves perfect classification
($F_1 = 1.0$) with median detection time $54.7\,\mu\mathrm{s}$
per program.

\paragraph{Future work.}
We plan to extend the framework to handle Python's descriptor
protocol (a fourth MOP-related mechanism), multiprocessing effects
(the \texttt{concurrent.futures} API), and signal handlers.  We
also plan to formalize the full soundness proof in a proof
assistant, using \JuGeo{}'s own framework to bootstrap the
verification.


%% ======================================================================
%% Bibliography
%% ======================================================================

\begin{thebibliography}{20}

\bibitem{jugeo-foundations}
H.~Young.
\newblock Judgment Geometry: A Sheaf-Theoretic Foundation for
  Multi-Source Program Verification.
\newblock \emph{Working paper}, 2025.

\bibitem{fstar-effects}
N.~Swamy, C.~Hri\c{t}cu, C.~Keller, A.~Rastogi, A.~Delignat-Lavaud,
  S.~Forest, K.~Bhargavan, C.~Fournet, P.-Y.~Strub, M.~Kohlweiss,
  J.-K.~Zinzindohou\'{e}, and S.~Zanella-B\'{e}guelin.
\newblock Dependent types and multi-monadic effects in {F*}.
\newblock In \emph{POPL}, pages 256--270, 2016.

\bibitem{dijkstra-monads}
D.~Ahman, C.~Hri\c{t}cu, K.~Maillard, G.~Mart\'{\i}nez,
  G.~Plotkin, J.~Protzenko, A.~Rastogi, and N.~Swamy.
\newblock Dijkstra monads for free.
\newblock In \emph{POPL}, pages 515--529, 2017.

\bibitem{steel-2021}
A.~Fromherz, A.~Rastogi, N.~Swamy, S.~Gibson, G.~Mart\'{\i}nez,
  D.~Merigoux, and T.~Ramananandro.
\newblock Steel: Proof-oriented programming in a dependently typed
  concurrent separation logic.
\newblock In \emph{ICFP}, 2021.

\bibitem{plotkin-power-2003}
G.~Plotkin and J.~Power.
\newblock Algebraic operations and generic effects.
\newblock \emph{Applied Categorical Structures}, 11(1):69--94, 2003.

\bibitem{bauer-pretnar-2015}
A.~Bauer and M.~Pretnar.
\newblock Programming with algebraic effects and handlers.
\newblock \emph{Journal of Logical and Algebraic Methods in
  Programming}, 84(1):108--123, 2015.

\bibitem{koka}
D.~Leijen.
\newblock Type directed compilation of row-typed algebraic effects.
\newblock In \emph{POPL}, pages 486--499, 2017.

\bibitem{frank}
L.~Convent, S.~Lindley, C.~McBride, and C.~McLaughlin.
\newblock Doo bee doo bee doo.
\newblock \emph{Journal of Functional Programming}, 30:e9, 2020.

\bibitem{iris}
R.~Jung, R.~Krebbers, J.-H.~Jourdan, A.~Bizjak, L.~Birkedal,
  and D.~Dreyer.
\newblock Iris from the ground up: A modular foundation for
  higher-order concurrent separation logic.
\newblock \emph{Journal of Functional Programming}, 28:e20, 2018.

\bibitem{crosshair}
P.~Curran.
\newblock {CrossHair}: Symbolic execution for Python contracts.
\newblock \url{https://github.com/pschanely/CrossHair}, 2020.

\bibitem{nagini}
M.~Eilers and P.~M\"{u}ller.
\newblock Nagini: A static verifier for {Python}.
\newblock In \emph{CAV}, pages 596--603, 2018.

\bibitem{deal}
G.~Grigorev.
\newblock Deal: Design by contract for {Python}.
\newblock \url{https://github.com/life4/deal}, 2019.

\bibitem{mypy}
J.~Lehtosalo et~al.
\newblock mypy: Optional static typing for {Python}.
\newblock \url{http://mypy-lang.org/}, 2012.

\bibitem{pyright}
E.~Traut.
\newblock Pyright: Static type checker for {Python}.
\newblock \url{https://github.com/microsoft/pyright}, 2019.

\bibitem{lean4}
L.~de~Moura and S.~Ullrich.
\newblock The {Lean~4} theorem prover and programming language.
\newblock In \emph{CADE}, pages 625--635, 2021.

\bibitem{dafny}
K.R.M.~Leino.
\newblock Dafny: An automatic program verifier for functional correctness.
\newblock In \emph{LPAR}, pages 348--370, 2010.

\end{thebibliography}

\end{document}
